% !TeX root = ../papers/nc-p3.tex
\section{単一二元辺を持つ全有限窓の直線複体と階数層}
\label{sec:binary-linear-complex}
\begin{theorem}\label{thm:binary-linear-complex}
一辺の二元商に向きと印を付けた有限点Frobenius窓は次の一つの既約なパラメータ空間の開部分で尽くされる。
$L:P_3\to P_5$ と $[\omega]\in\mathbb P(\bigwedge^2 P_3)$ に対し
\[
 F_L(q\wedge r)=qL(r)-rL(q)\in P_8,\qquad F_L(\omega)=0
\]
を課す。$F_L$ の階数は5とし、$S=(\operatorname{im}F_L)^\perp\subset P_8^*$ とする。
窓を $Q,U,Q,S$ 順に、$Q=P_3$, $U=P_2\oplus k$ として
\[
  w(q,(p,a),q',D)=D\big((pq+aL(q))q'\big)
\]
で与える。これに有限点グラフと有限Frobenius検査を課すだけで全許容値となり、全て非正則である。
第四文字を加える自由度 $L\mapsto sL+M_p$ と同時の二元座標変換を除く。
ここで $M_p$ は $p$ による乗法である。

非退化な $\omega$ なら窓の反対階数は10で、この層は稠密である。
$\omega=q\wedge r$ が分解するとき、$e=\deg\gcd(q,r)$ とする。
$e=0$ は有限グラフを持たず、$e=1$ は反対階数9、$e=2$ は反対階数8である。
後二層とも有限Frobenius窓が実在し、$e=2$ はこの向きの両対称二元窓に一致する。
従って非正則な有限窓の中に、階数10で安定・最小周期4・二十点全て被約の
既約な13次元の添字保持同型類の族がある。この族の一般点で二元商の印は一意に戻る。
各有限非両対称点では、商を取る位置は四位置のうちただ一つで、最小周期は常に4である。
\end{theorem}
\begin{proof}[全パラメータと既約性]
前定理で $Q$ の両位置は同定され、第四文字による縮約は $Q$ の交換で不変である。
二元乗法を $L$ で拡張すれば、その不変性は $D$ が $F_L$ の像を消すこととなる。
実際、三次交叉 $Q,U,Q$ は $P_2$ への制限で $P_8^*$ に入り、両辺のHankel条件から
四文字目の値まで一意に決まる。従ってこの交叉はちょうど $(\operatorname{im}F_L)^\perp$ で、
Frobenius条件により階数5である。逆に上の式は両辺で指定の関係を実現し、
開な有限Frobenius条件とグラフの有限性を通常の有限行列・入射検査で決定できる。
$L$ の第四文字の変更と二元座標を除けば、印を保つ同型条件も必要十分である。
これは圧縮辺の階数一五次曲線が二つの座標を共通の変換まで決め、$S$ は三次交叉で内在的に決まるためである。

非退化な $\omega$ では $L\mapsto F_L(\omega)$ は $P_8$ に全射である。
実際非退化な四行交代行列で単項式基底を変換した四つの三次式と任意の五次式の積を足す写像である。
従って $\mathbb P^5$ の開部分上で、入射は階数15のベクトル束である。
分解するときの像は $qP_5+rP_5=dP_{8-e}$ ($d=\gcd(q,r)$) で、核は $15+e$ 次元。
$[d]$ と基点自由な二平面 $\langle q/d,r/d\rangle\subset P_{3-e}$ は計 $4-e$ 次元を動く。
よって $e=0,1,2$ の入射は各19次元である。全入射は29次元の滑らかな空間で九式により定まり、
各成分は20次元以上なので、これらは全て既約な非退化部分の閉包に入る。

$e=0$ では $L(q)=cq,L(r)=cr$ ($c\in P_2$) となり $(-c,1)\in U$ の評価核が二次元以上となる。
$e=2$ なら $\langle q,r\rangle=dP_1$ で、$L(dv)=cv$ ($v\in P_1$, $c\in P_4$)。
従って前定理群で用いた三差分は $d$ の係数による定数関係を持つ。
純関係がないためこれらは基点自由であり、六次の鉛筆を張る。
その完全交叉の十次汎関数を取り $GL$ がHankelになることから
定理\ref{thm:binary-apolar-rank-test}の両対称窓に戻る。
逆に全両対称二元点では、核二次式 $d$ に対する同じ式から $\omega$ はこの分解型である。
$e=1$ の有限例は定理\ref{thm:binary-reverse-orientation}の $d=d_1d_2$, $c=d_2c_3$ の有限例から得る。
指定の二平面を $d_1\langle g_t,h_t\rangle$ として
$(g_0,h_0)=d_2(s,t')$ から互いに素な二次式対へ動かす（$s,t'$ は二元変数）。
$L_t(d_1g_t)=c_3g_t$, $L_t(d_1h_t)=c_3h_t$ とし、補二文字の値も延ばせば所望の $L_0$ に特殊化する。
階数5、有限性、Frobenius性は特殊点に成立する開条件である。
稠密な非退化 $\omega$ の側でもこれらが成り立ち、一般二元点からの開性により点は一般に被約である。
\end{proof}
\begin{proof}[射影直線上の二平面束と反対階数]
$v(t),\lambda(t)\in Q^*$ をそれぞれ $q(t),L(q)(t)$ の評価とする。
これらは純関係の不在により各点で独立で、生成する二平面は
$\mathcal E=\mathcal O(-3)\oplus\mathcal O(-5)\subset Q^*\otimes\mathcal O$ を与える。
そのPlücker座標系 $H=\operatorname{im}F_L$ は $\mathbb P^1$ を $\mathbb P^4$ へ埋め込む。
実際長さ2を分離できなければ $S=H^\perp$ に長さ2以下の因子に支持される八次汎関数がある。
その $H_{5,3}$ の核は二次元以上となり、$Q,S$ 辺の有限性に反する。

反対切断を $\operatorname{Sym}^2 Q$ からの写像とする。その核の二次式 $A$ の
対角制限 $f=A(v,v)\in P_6$ は $fP_2\subset H$ を満たす。
$f\ne0$ ならその長さ6の零因子で $H$ の像は高々二次元で、その埋込み像は直線に含まれる。
Plücker二次式はこの直線全体で零である。Grassmann多様体の直線は固定した平面内の直線の鉛筆だから、
長さ6の零因子に沿い $v(t)$ がその固定平面に入ってしまう。三次正規曲線の平面切断は長さ3で矛盾する。
よって $A(v,v)=0$。第四文字の条件は $A(v,\lambda)\in H$ で、
これは適当な交代行列を $A$ に足した行列 $B$ で $B(v,\lambda)=0$ となることに等しい。
同時に $B(v,v)=0$。$B(v,-)$ は任意の三次係数の四成分を一意に表せるので、
反対核の次元は $h^0(\mathcal E^\perp(3))-1$ である。引いた一次元は交代行列 $\omega$ の核線である。

非退化なら $\omega$ で $\mathcal E^\perp\simeq\mathcal E$ だから反対核は零である。
ここで反対核は対称二乗十次元内で取っている。
分解型なら $\omega(v,-)/d$ は次数 $3-e$ の零点を持たない切断列となる。
従って $\mathcal E^\perp$ は次数 $-(3-e)$ の線部分束を持ち、商の次数は $-(5+e)$。
この拡大は $\mathbb P^1$ 上で分裂し、反対核の次元は $e$ となる。これで全階数を得る。
\end{proof}
\begin{proof}[印を忘れた像と周期]
非退化部分の20次元から第四文字の自由度4と $\operatorname{PGL}_2$ の3を除き、一般の印付き像は
高々13次元である（一般の二元四次成分を持つ $[\omega]$ の射影安定群は有限）。
その閉包は12次元の両対称二元軌跡を含み、稠密部の反対階数は10であるから、
印を忘れた像の次元も13である。一般の印の一意性には命題\ref{prop:binary-quotient-rank-two-curve}と
定理\ref{thm:binary-incidence-codimension}の残余平面六次を使える。
\path{computations/verify_binary_linear_complex.py}の非退化標本ではその六次が滑らかである。
階数一Hankel点 $h(t)$ が残余六次の一般点の核になることはない。
その場合の二次式は $t$ で消え、残る一次因子は追加行と二つの条件を満たす。
純関係がなければその一次因子はあるとしても一意である
（二条件がともに恒等零なら追加行は $t$ の評価の倍数になる）。
従ってこの型の曲線は有理で、滑らかな平面六次を支配できない。
よって残余成分はその六次と双有理で、関係六平面の唯一の有理正規五次曲線が指定の印を復元する。

両対称でない窓では、別の向きで全三次文字を $U$ または $S$ とする商があれば、
前定理によりもう一方の反対交換も対称となってしまう。
同じ $Q$ を全三次文字、$S$ を圧縮側にする商も、既存の $U$ を圧縮する辺と
$Q$ を挟んで隣接するため定理\ref{thm:binary-two-edge-rigidity}に反する。
従って圧縮する位置は一意であり、二移動では保持されない。四移動と合わせて最小周期を得る。
非正則有限窓は定理\ref{thm:all-finite-windows-semistable}から安定である。
\end{proof}
